\documentclass[11pt,a4paper]{article}
\usepackage[T1]{fontenc}
\usepackage[utf8]{inputenc}
\usepackage{lmodern}
\usepackage[margin=25mm]{geometry}
\usepackage{amsmath,amsthm,amssymb,mathtools}
\usepackage{microtype}
\usepackage{booktabs,longtable,array,tabularx}
\usepackage{enumitem}
\usepackage{xcolor}
\usepackage{listings}
\usepackage[colorlinks=true,linkcolor=blue!40!black,citecolor=blue!40!black,urlcolor=blue!50!black]{hyperref}
\hypersetup{pdftitle={Decomposing algebraic numbers into sums and products of minimal degree}}

\lstset{basicstyle=\ttfamily\small,columns=fullflexible,keepspaces=true,breaklines=true,
  frame=single,framerule=0.3pt,rulecolor=\color{black!30},xleftmargin=4pt,xrightmargin=4pt,
  aboveskip=6pt,belowskip=6pt,showstringspaces=false}

\newtheorem{theorem}{Theorem}[section]
\newtheorem{proposition}[theorem]{Proposition}
\newtheorem{lemma}[theorem]{Lemma}
\newtheorem{corollary}[theorem]{Corollary}
\theoremstyle{definition}
\newtheorem{definition}[theorem]{Definition}
\newtheorem{example}[theorem]{Example}
\newtheorem{problem}[theorem]{Further problem}
\theoremstyle{remark}
\newtheorem{remark}[theorem]{Remark}

\newcommand{\Q}{\mathbb{Q}}
\newcommand{\Z}{\mathbb{Z}}
\newcommand{\C}{\mathbb{C}}
\newcommand{\R}{\mathbb{R}}
\newcommand{\F}{\mathbb{F}}
\newcommand{\Qbar}{\overline{\Q}}
\newcommand{\Gal}{\operatorname{Gal}}
\newcommand{\Tr}{\operatorname{Tr}}
\newcommand{\Nm}{\operatorname{N}}
\newcommand{\Res}{\operatorname{Res}}
\newcommand{\Stab}{\operatorname{Stab}}
\newcommand{\lcm}{\operatorname{lcm}}
\newcommand{\rk}{\operatorname{rank}}
\newcommand{\dg}{\deg}
\newcommand{\Dp}{D_{+}}
\newcommand{\Dt}{D_{\times}}
\newcommand{\Dtt}{D_{\times,2}}
\newcommand{\Dpp}{D_{+,2}}
\newcommand{\code}[1]{\texttt{#1}}
\newcommand{\Root}{\code{Root}}

\title{Decomposing algebraic numbers into sums and products of minimal degree\\[4pt]
\large A unified account of the theory, with complete algorithms for sums and for two-factor products,
and an implementation in the Wolfram Language and in Python}
\author{Prepared for Vladimir Reshetnikov\\ \small in answer to Mathematica StackExchange question 105933}
\date{September 2026}

\begin{document}
\maketitle

\begin{abstract}
Given an algebraic number $a$ of degree $n$, we ask for a representation of $a$ as a sum, or as a
product, of algebraic numbers whose largest degree is as small as possible.  This article unifies nine
independent technical reports written on the question and adds several results.  We prove that the
additive problem has a complete finite algorithm (trace descent to the splitting field, then rational
linear algebra on fixed fields of subgroups of the Galois group), that the two-factor multiplicative
problem has a complete finite algorithm (norm descent with an explicit radical exponent), and that the
arbitrary-length multiplicative problem is different: optimal factors may lie outside the
splitting field, two-factor optima are not many-factor optima, and the methods developed here do not
give a complete algorithm for arbitrary products.  We
prove rigorous lower bounds valid against any finite number of components, including a computable bound
from Frobenius cycle types, and we settle the two examples of the question: both have globally optimal
maximum degree three.  We also answer a further question: a \Root{} can be a flat sum of \Root{}s of
smaller degree, even with Gaussian rational coefficients, while no sequence of binary additive or
multiplicative splittings into two \Root{}s of smaller degree exists; the number $\sqrt2+\sqrt3+\sqrt6$
is such an example.  The accompanying Wolfram Language package and its Python counterpart implement
the principal algorithms with an engine based on numerical resolvents for the Galois group and a trace-form basis of the
splitting field.  This avoids the slow number-field conversions proposed in the reports;
exact identity verification and the certification of exhaustive searches are discussed separately.
\end{abstract}

\tableofcontents

\section{Introduction}

The question \cite{question} asks, for a polynomial \Root{} object such as
\begin{equation}\label{eq:examples}
\begin{aligned}
 \alpha_\times&=\Root[-1 - \# + 3 \#^3 - \#^4 + \#^5 - 3 \#^6 + 2 \#^7 + \#^9 \&, 1],\\
 \alpha_+&=\Root[8 - 4 \# + 24 \#^2 - 15 \#^3 + 3 \#^5 + 6 \#^6 + \#^9 \&, 1],
\end{aligned}
\end{equation}
for a representation as a product, respectively a sum, of \Root{} objects of the smallest possible
maximum degree.  The expected answers are
\begin{equation}\label{eq:answers}
\begin{aligned}
 \alpha_\times&=\Root[1+\#+\#^3\&,1]\cdot\Root[1-\#+\#^3\&,1],\\
 \alpha_+&=\Root[1+\#+\#^3\&,1]+\Root[1-\#+\#^3\&,1].
\end{aligned}
\end{equation}
The constructive answer on the site \cite{question} guesses the shape of the answer (two cubics), forms
the characteristic polynomial of a Kronecker product or Kronecker sum of companion matrices with
symbolic coefficients, and solves the resulting coefficient equations.  This finds the decomposition once
its shape is known, but it neither discovers the shape nor proves that no better one exists.

Nine technical reports (\cite{r1,r2,r3,r4,r5,r6,r7,r8,r9}, reproduced verbatim in the repository)
were written on the question.  They agree on the core theory and differ in emphasis, examples and
implementation.  None of them executed its Wolfram Language code, and the number-field construction they
all propose turns out to be impractical in Mathematica (Section~\ref{sec:impl}).  This article gives one
coherent account of the mathematics, proves everything that is used, records the examples and
counterexamples that the reports contributed, adds a treatment of binary versus flat decompositions and
of Gaussian rational coefficients, and describes an implementation that has actually been run.

\subsection{Summary of results}

Write $\dg(b)=[\Q(b):\Q]$ for the absolute degree of an algebraic number $b$, and define
\begin{align*}
 \Dp(a)&=\min\Bigl\{\max_i\dg(b_i): a=b_1+\dots+b_r,\ r\ge1\Bigr\},\\
 \Dt(a)&=\min\Bigl\{\max_i\dg(b_i): a=b_1\cdots b_r,\ r\ge1\Bigr\},
\end{align*}
the components ranging over all algebraic numbers.  The main results are the following.
\begin{enumerate}[label=(\roman*)]
\item \emph{Lower bounds} (Section~\ref{sec:bounds}).  Every prime divisor of $\dg(a)$ is at most
  $\Dp(a)$ and at most $\Dt(a)$; the exponent of the Galois group $G$ of the normal closure of $\Q(a)$
  divides $\lcm(1,\dots,d)$ for $d=\Dp(a)$ or $d=\Dt(a)$; and orders of Frobenius elements, read off from
  factorization patterns modulo primes, give a rigorous lower bound without computing $G$.
\item \emph{Sums} (Section~\ref{sec:sums}).  $\Dp(a)\le d$ if and only if $a$ lies in the $\Q$-linear
  span of the subfields of degree at most $d$ of the splitting field $L$ of $a$.  Equivalently, $a$ lies
  in the span of the subset sums of its conjugates of degree at most $d$.  This is a complete finite
  algorithm.
\item \emph{Two-factor products} (Section~\ref{sec:products}).  $a=bc$ with $\dg b,\dg c\le d$ if and
  only if there are subfields $E,F\subseteq L$ and an integer $t\ge1$ with $t[E:\Q]\le d$, $t[F:\Q]\le d$
  and $E\cap a^tF\ne\{0\}$.  The exponent $t$ cannot be dropped: for $a=\sqrt{(1+\sqrt2)(1+\sqrt3)}$
  the optimal factors $\sqrt{1+\sqrt2}$ and $\sqrt{1+\sqrt3}$ lie outside $L$, and no elements of $L$
  of degree below $8$ multiply or add up to $a$.
\item \emph{Many factors} (Section~\ref{sec:many}).  $\Dt$ can be smaller than the two-factor
  minimum: $(1+\sqrt2)(1+\sqrt3)(1+\sqrt5)$ has $\Dt=2$ but two-factor minimum $4$.
  Products inside a fixed field are a multiplicative, not a linear, problem; arithmetic obstructions
  exist beyond Galois theory ($\Dp(1+\sqrt2+\sqrt3)=2$ but $\Dt(1+\sqrt2+\sqrt3)=4$).  We give the
  necessary norm condition, a sufficient radical condition, and an exact rank-one tensor criterion for
  linearly disjoint families; a complete many-factor decision procedure remains outside this work.
\item \emph{Binary versus flat} (Section~\ref{sec:binary}).  A decomposition into two components of
  smaller degree, iterated, is strictly weaker than a flat decomposition: $\sqrt2+\sqrt3+\sqrt6$ is a sum
  of three quadratic numbers but is neither a sum nor a product of two algebraic numbers of degree below
  four, even if Gaussian rational coefficients are allowed at every binary step.  With Gaussian
  coefficients the flat additive problem is again linear algebra, in $L(i)$.
\item \emph{The examples} (Section~\ref{sec:examples}).  Both numbers in \eqref{eq:examples} have
  globally optimal maximum degree $3$, attained by \eqref{eq:answers}.  Moreover $\Dp(\alpha_\times)=6$,
  attained by two sextics that do not lie in $\Q(\alpha_\times)$, whereas every additive decomposition
  of $\alpha_\times$ inside $\Q(\alpha_\times)$ has maximum degree $9$.
\item \emph{Implementation} (Section~\ref{sec:impl}).  The Galois group is computed by numerical
  resolvents, the splitting field is represented in a tower basis with coordinates obtained from traces,
  and the additive and two-factor criteria become rational linear algebra once the field data are certified.
  The many-factor search also needs multiplicative tests.  A fast path works inside
  $\Q(a)$ without the Galois group.  The original measurements answered the two examples in about a
  second in Wolfram (fast path), certifying global optimality; the Python port (python-flint, Arb ball
  arithmetic) constructed the full splitting-field data in three to five seconds, against about
  twenty seconds in the Wolfram kernel.  Section~\ref{sec:impl} distinguishes these historical timings
  from validation of the revised code.
\end{enumerate}

\subsection{Conventions}

All fields are subfields of $\C$, so every algebraic extension is separable.  A number field is a
finite extension of $\Q$.  For a finite Galois extension $L/\Q$ with group $G=\Gal(L/\Q)$, subfields of
$L$ correspond to subgroups $H\le G$ via $H\mapsto L^H$, and $[L^H:\Q]=[G:H]$.  The degree of an element
$b\in L$ is the index of its stabilizer, equivalently the size of its $G$-orbit.  We write $\Res_y(f,g)$
for the resultant with respect to $y$, and $\Tr_{M/L}$, $\Nm_{M/L}$ for relative trace and norm.

Mathematica's \code{Root[f,k]} indexes real roots first, in increasing order, and places non-real
conjugate pairs consecutively for rational polynomials.  Its documentation warns that the order of
non-real roots can depend on the root-isolation method \cite{WolframRoot}.  The Python code uses its
own explicit order for those pairs: increasing real part, then increasing absolute imaginary part,
with the negative imaginary part first.  Thus real-root indices agree, including the two original
examples, but an arbitrary non-real branch must be matched by an isolating region when moving between
implementations; a printed index alone is not a universal interchange format.

\section{The problem}\label{sec:problem}

\begin{definition}\label{def:D}
For $a\in\Qbar$ put
\begin{align*}
 \Dp(a)&=\min\{\max_i\dg(b_i): r\ge1,\ b_i\in\Qbar,\ a=b_1+\dots+b_r\},\\
 \Dt(a)&=\min\{\max_i\dg(b_i): r\ge1,\ b_i\in\Qbar^\times,\ a=b_1\cdots b_r\}\quad(a\ne0),
\end{align*}
and $\Dt(0)=1$.  Let $\Dpp(a)$ and $\Dtt(a)$ be the same minima restricted to $r\le2$.  For a field
$K\ni a$ let $\Dp^K(a)$, $\Dt^K(a)$ be the minima with all components in $K$.
\end{definition}

The one-component representation shows that all these quantities are at most $\dg(a)$, and
$\Dt\le\Dtt$, $\Dp\le\Dpp$, $\Dp\le\Dp^K$, $\Dt\le\Dt^K$.  Several elementary remarks fix the
conventions.

\begin{enumerate}[label=(\alph*)]
\item Degrees are absolute.  A displayed \Root{} may be a root of a reducible polynomial, and a
  polynomial of small degree with algebraic coefficients says nothing about absolute degree: $x-a$ has
  degree one over $\Q(a)$.  Degrees are therefore measured with the rational minimal polynomial.
\item Rational scalars are free: for $q\in\Q^\times$, $\Q(qb)=\Q(b)$.  A rational linear combination of
  low-degree numbers is a permitted sum, and a rational multiple of a factor may be moved to another
  factor.  Rational numbers have degree one.  Inverses and negatives do not change degrees.
\item Components are allowed to be complex even when $a$ is real; requiring real components is a
  different problem (Remark~\ref{rem:real}).
\item The number of components, the coefficient heights and the printed size of the answer are not
  minimized; they are secondary objectives.  A flat sum $b_1+\dots+b_r$ is not a nested expression such
  as $b_1+b_2b_3$; Section~\ref{sec:binary} compares the two.
\item Three ambient domains matter: $K=\Q(a)\subseteq L\subseteq\Qbar$, where $L$ is the splitting
  field of the minimal polynomial of $a$.  For sums, $L$ is always large enough; for products not even
  $L$ is (Section~\ref{sec:products}).  Restricting to $K$ can change both answers.
\end{enumerate}

\section{Composed sums and products}\label{sec:resultants}

Let $f,g\in\Q[x]$ be monic of degrees $m,k$ with roots $u_1,\dots,u_m$ and $v_1,\dots,v_k$ (with
multiplicity).  Their \emph{composed sum} and \emph{composed product} are
\begin{equation}\label{eq:composed}
\begin{aligned}
 (f\oplus g)(x)&=\Res_y\bigl(f(y),g(x-y)\bigr)=\prod_{i,j}(x-u_i-v_j),\\
 (f\otimes g)(x)&=\Res_y\bigl(f(y),\,y^k g(x/y)\bigr)=\prod_{i,j}(x-u_iv_j).
\end{aligned}
\end{equation}
Both identities follow from $\Res_y(f,h)=\prod_i h(u_i)$ for monic $f$; the second argument of the
product resultant is the polynomial $\sum_j g_jx^jy^{k-j}$, which handles a zero root correctly.
Equivalently, $f\oplus g$ and $f\otimes g$ are the characteristic polynomials of $C_f\otimes I+I\otimes
C_g$ and $C_f\otimes C_g$, where $C_f,C_g$ are companion matrices; this is the construction used in the
answer on the site.  Fast algorithms for composed sums and products exist \cite{BFSS}; here only small
degrees occur.

\begin{proposition}\label{prop:divides}
Let $p$ be the minimal polynomial of $a$.  Some pair of roots $(u_i,v_j)$ satisfies $a=u_i+v_j$
(respectively $a=u_iv_j$) if and only if $p$ divides $f\oplus g$ (respectively $f\otimes g$).
\end{proposition}

\begin{proof}
If $a=u_i+v_j$ then $a$ is a root of $f\oplus g$, and the irreducible $p$ divides it.  Conversely, if
$p\mid f\oplus g$ then $a$ is a root of the product in \eqref{eq:composed}, so some linear factor
vanishes at $a$.  The product case is identical.
\end{proof}

Divisibility cannot be replaced by equality: $(x^2-2)\oplus(x^2-2)=x^2(x^2-8)$ and
$(x^2-2)\otimes(x^2-3)=(x^2-6)^2$, while $\sqrt2+\sqrt2$ and $\sqrt2\cdot\sqrt3$ have minimal
polynomials $x^2-8$ and $x^2-6$.  Nor need $\dg(a)$ divide $mk$: with $u=\sqrt[3]{2}$ and
$\omega=e^{2\pi i/3}$ the number $u-\omega u$ has minimal polynomial $x^6+108$, of degree $6\nmid9$.
The correct general statement is $\dg(a)\le[\Q(b_1,\dots,b_r):\Q]\le\prod\dg(b_i)$.

For the two cubics of \eqref{eq:answers} one computes
\begin{equation}\label{eq:resultants}
\begin{aligned}
 (x^3+x+1)\oplus(x^3-x+1)&=x^9+6x^6+3x^5-15x^3+24x^2-4x+8,\\
 (x^3+x+1)\otimes(x^3-x+1)&=x^9+2x^7-3x^6+x^5-x^4+3x^3-x-1,
\end{aligned}
\end{equation}
exactly the polynomials in \eqref{eq:examples}.

\subsection{Inverse composition: fix one polynomial}

Guessing both polynomials with unknown coefficients leads to a nonlinear system whose solutions over
$\C$ may have algebraic coefficients (which do not certify absolute degrees) and whose normalizations
need justification (rescaling $u\mapsto qu$, $v\mapsto v/q$ with $q\in\Q^\times$ preserves a product,
but making a constant term equal to $1$ requires a rational $m$th root).  A cheaper exact device fixes
one candidate polynomial $f$ and recovers the partner: if $a=u+v$ with $f(u)=0$ then $v$ is a root of
$\Res_y(f(y),p(x+y))$, and if $a=uv$ then $v$ is a root of $\Res_y(f(y),p(xy))$ (for $f(0)\ne0$).
Factoring these resultants over $\Q$ and keeping the factors of degree at most $d$ gives a complete
search for a fixed first polynomial.  For $f=x^3+x+1$ and both examples the resultant factors as
$(x^3-x+1)^3h_{18}(x)$ with $h_{18}$ irreducible \cite{r2,r4}.  Enumerating first polynomials of
bounded degree and coefficient height, including non-monic primitive polynomials (otherwise
$\tfrac12$ would be missed), gives the \emph{bounded dictionary search}: complete within its box, a
positive semi-decision procedure in general, never a proof of impossibility.  Even the smallest box,
cubics with coefficients in $\{-1,0,1\}$, discovers both cubics of \eqref{eq:answers}.

\section{The examples: identities, degrees, branches}\label{sec:examples}

Let $u$ and $v$ be the real roots of $f=x^3+x+1$ and $g=x^3-x+1$.  Both cubics are irreducible (no
rational roots) with discriminants $-31$ and $-23$, so each has exactly one real root and Galois group
$S_3$.  Write $P_\times$ and $P_+$ for the two polynomials in \eqref{eq:examples}.

\begin{proposition}\label{prop:examples}
$P_\times$ and $P_+$ are irreducible, each has exactly one real root, and $\alpha_\times=uv$,
$\alpha_+=u+v$.  Moreover $\Q(uv)=\Q(u+v)=\Q(u,v)$ has degree nine, and
\begin{equation}\label{eq:recover}
 u=\frac{\alpha_\times^3-\alpha_\times^2-1}{\alpha_\times^4+\alpha_\times^2-\alpha_\times+1},\quad
 v=\frac{\alpha_\times}{u};\qquad
 u=\frac{3\alpha_+^5-5\alpha_+^3-3\alpha_+^2+2\alpha_+-4}{6\alpha_+^4-6\alpha_++4},\quad v=\alpha_+-u.
\end{equation}
Also $u=-(\alpha_\times^7+2\alpha_\times^5-2\alpha_\times^4+1)/2$ and
$v=-(\alpha_\times^7+2\alpha_\times^5-2\alpha_\times^4+2\alpha_\times^3+2\alpha_\times-1)/2$.
\end{proposition}

\begin{proof}
The splitting fields $L_f,L_g$ of the cubics are $S_3$-extensions whose unique quadratic subfields
$\Q(\sqrt{-31})$ and $\Q(\sqrt{-23})$ differ.  A nontrivial common Galois quotient of two copies of
$S_3$ is $C_2$ or $S_3$, and either would force a common quadratic subfield; hence $L_f\cap L_g=\Q$ and
$\Gal(L_fL_g/\Q)=S_3\times S_3$, acting transitively on the nine ordered pairs $(u_i,v_j)$ of roots.
If two pairs had the same sum, $u_i-u_k=v_l-v_j\in L_f\cap L_g=\Q$; a nonzero rational difference of
two conjugates is impossible, since iterating an automorphism mapping $u_k\mapsto u_i$ would produce
arbitrarily large differences.  If two pairs had the same product, $u_i/u_k\in\Q^\times$ would be a root
of unity, hence $\pm1$, and $-1$ is excluded because $f(x)+f(-x)=2\ne0$.  So the nine sums and the nine
products are distinct, the resultants \eqref{eq:resultants} are the minimal polynomials, and both are
irreducible of degree nine.  A real conjugate of $uv$ or $u+v$ is fixed by complex conjugation, which
permutes the pairs; the only pair fixed by conjugation is $(u,v)$, so each polynomial has exactly one
real root, namely $uv$, respectively $u+v$; by the root ordering this is \code{Root[\ldots,1]}.

The recovery formulas are verified by exact polynomial arithmetic modulo $P_\times$ and $P_+$: the
denominators are coprime to the respective polynomial, and substituting the right-hand sides into
$f$ and $g$ gives zero modulo $P_\times$, respectively $P_+$.  The real-branch identification follows
from the uniqueness of the real roots of the two cubics.  These checks show $u,v\in\Q(\alpha_\times)$ and
$u,v\in\Q(\alpha_+)$; since $[\Q(u,v):\Q]=9$, all three fields coincide.
\end{proof}

\section{Lower bounds valid for any number of components}\label{sec:bounds}

\begin{theorem}[normal closure obstruction]\label{thm:bound}
Let $n=\dg(a)$, let $L$ be the splitting field of the minimal polynomial of $a$ and $G=\Gal(L/\Q)$.
Suppose $a$ lies in the field generated by finitely many algebraic numbers $b_1,\dots,b_r$ of degrees at
most $d$ (in particular, if $a$ is a sum or a product of them).  Then
\begin{enumerate}[label=(\roman*)]
\item every prime divisor of $|G|$, hence of $n$, is at most $d$;
\item $\exp(G)$ divides $\lcm(1,2,\dots,d)$; in particular $d\ge$ the largest prime power dividing $\exp(G)$;
\item if $d\le4$ then $G$ is solvable.
\end{enumerate}
Consequently $\Dp(a)$ and $\Dt(a)$ are at least the largest prime divisor of $n$ and at least the
largest prime power dividing $\exp(G)$.
\end{theorem}

\begin{proof}
Let $M_i$ be the normal closure of $\Q(b_i)$ and $M=M_1\cdots M_r$.  $\Gal(M_i/\Q)$ acts faithfully on
the $m_i\le d$ conjugates of $b_i$, so it embeds in $S_{m_i}$, and restriction embeds $\Gal(M/\Q)$ into
$\prod_i\Gal(M_i/\Q)\le\prod_iS_{m_i}$.  Every prime divisor of $|\prod S_{m_i}|=\prod m_i!$ is at most
$d$, the exponent of $\prod S_{m_i}$ divides $\lcm(1,\dots,d)$ (orders of permutations are lcm's of cycle
lengths), and $S_{m}$ is solvable for $m\le4$.  These properties pass to subgroups and to quotients.
Since $M/\Q$ is normal and contains $a$, it contains all conjugates of $a$, hence $L$, and restriction
maps $\Gal(M/\Q)$ onto $G$.  Finally $n=[\Q(a):\Q]$ divides $|G|=[L:\Q]$.
\end{proof}

For the examples, $n=9$ gives the bound $3$, which \eqref{eq:answers} attains; so
$\Dt(\alpha_\times)=\Dp(\alpha_+)=3$ globally, against any finite number of complex components.  The
exponent bound is genuinely stronger than the prime bound: for $\zeta_5$ ($G=C_4$) and for
$\sqrt{2+\sqrt2}$ ($G=C_4$) it gives $\Dp=\Dt=4$ although the prime bound is $2$; for $2^{1/4}$
($G=D_4$, exponent $4$) it gives $4$; for a root of $x^4-x-1$ ($G=S_4$, exponent $12$) it gives $4$.

\begin{proposition}[Frobenius bound]\label{prop:frobenius}
Let $P\in\Z[x]$ be the primitive minimal polynomial of $a$ and let $q$ be a prime dividing neither the
leading coefficient nor the discriminant of $P$.  If $P$ factors modulo $q$ into irreducible factors of
degrees $e_1,\dots,e_s$, then $G$ contains an element of order $\lcm(e_1,\dots,e_s)$, so this number
divides $\exp(G)$.  Hence $\Dp(a),\Dt(a)\ge$ the largest prime power dividing
$\lcm_q\lcm(e_1,\dots,e_s)$, the outer lcm over any finite set of such primes.
\end{proposition}

\begin{proof}
For unramified $q$ the Frobenius element at a prime of $L$ above $q$ acts on the roots with cycle type
$(e_1,\dots,e_s)$ (Dedekind's theorem, e.g.\ \cite[Th.~4.37]{Cohen}), hence has order
$\lcm(e_1,\dots,e_s)$.
\end{proof}

The Frobenius bound is cheap (a few factorizations modulo small primes) and is the bound the
implementations use before any Galois group is computed.  For $x^4-x-1$ the factorizations modulo $2$
and $7$ have patterns $(4)$ and $(1,3)$, giving exponent multiple $12$ and the bound $4$.  In general
the bound may be weaker than the true exponent bound; the full Galois data, when available, sharpens
it.

\begin{proposition}[length-dependent bound]\label{prop:length}
If $a=b_1\circ\dots\circ b_r$ (sums or products) with $\dg(b_i)\le d$, then $n\le d^r$.  In particular a
two-component representation needs $d\ge\lceil\sqrt n\rceil$.
\end{proposition}

\begin{proof}
$\Q(a)\subseteq\Q(b_1,\dots,b_r)$, whose degree is at most $\prod\dg(b_i)$.
\end{proof}

The bounds of this section are necessary conditions.  They are not sufficient even for sums:
$\alpha_\times$ has $\Dp=6$ (Example~\ref{ex:sextic}), while its degree-nine prime bound and its
$S_3\times S_3$ exponent bound are both $3$.  Nor are the bounds sufficient for products even when the
Galois group is elementary abelian (Theorem~\ref{thm:sign}).

\section{Sums: a complete finite algorithm}\label{sec:sums}

\subsection{Trace descent}

\begin{lemma}[trace projection]\label{lem:trace}
Let $L/\Q$ be a finite Galois extension, $a\in L$, and $a=b_1+\dots+b_r$ with arbitrary algebraic $b_i$.
Then there are $c_i\in L\cap\Q(b_i)$ with $a=c_1+\dots+c_r$; in particular $\dg(c_i)\le\dg(b_i)$.
\end{lemma}

\begin{proof}
Choose a finite Galois extension $M/\Q$ containing $L$ and all $b_i$; put $\Gamma=\Gal(M/\Q)$ and
$N=\Gal(M/L)$, a normal subgroup of $\Gamma$ since $L/\Q$ is Galois.  Define the $\Q$-linear projection
\[
 \pi(z)=\frac1{|N|}\sum_{\sigma\in N}\sigma(z)=\frac{\Tr_{M/L}(z)}{[M:L]},\qquad M\to L,
\]
which fixes $L$ pointwise.  Normality of $N$ gives $\gamma\pi(z)=\pi(\gamma z)$ for all
$\gamma\in\Gamma$: conjugation by $\gamma$ permutes $N$.  Hence every automorphism fixing $b_i$ fixes
$\pi(b_i)$, and by the Galois correspondence $\pi(b_i)\in\Q(b_i)$.  Put $c_i=\pi(b_i)\in L\cap\Q(b_i)$;
then $\sum c_i=\pi(a)=a$.
\end{proof}

Normality of $L$ is essential: averaging over $\Gal(M/\Q(a))$ for non-normal $\Q(a)$ is not
$\Gamma$-equivariant, and Example~\ref{ex:quartic} shows that restricting summands to $\Q(a)$ can
change the answer.

\subsection{The fixed-space criterion}

For $L$ Galois with group $G$ and $d\ge1$ define the $\Q$-vector subspace
\begin{equation}\label{eq:Wd}
 W_d(L)=\sum_{H\le G,\ [G:H]\le d}L^H\ \subseteq L,
\end{equation}
a sum of subspaces, not a compositum of fields.

\begin{theorem}[complete additive criterion]\label{thm:sum}
Let $a\in L$ with $L/\Q$ finite Galois.  Then $\Dp(a)\le d$ if and only if $a\in W_d(L)$.  In
particular $\Dp(a)=\min\{d: a\in W_d(L)\}$, and this is computable by finitely many rational
linear-algebra operations once $G$ and its subgroups are known.  If $a$ is not rational the minimum
is attained by a representation with at most $[L:\Q]$ nonzero summands.
\end{theorem}

\begin{proof}
If $a=\sum b_i$ with $\dg(b_i)\le d$, Lemma~\ref{lem:trace} gives $a=\sum c_i$ with $c_i\in L$ and
$\dg(c_i)\le d$; each $c_i$ lies in the subfield $\Q(c_i)=L^{H_i}$ with $[G:H_i]\le d$, so
$a\in W_d(L)$.  Conversely an element of $W_d(L)$ is a finite sum of elements of fields of degree at
most $d$, each of which has degree at most $d$.  Choosing a basis of $W_d(L)$ from the union of bases of
the fields and absorbing the rational coefficients gives at most $\dim W_d(L)\le[L:\Q]$ summands.
The minimum is attained at some $d\le\dg(a)$ because $\Q(a)\subseteq L$ is itself one of the fields.
\end{proof}

Only the subgroups $H$ that are minimal among those of index at most $d$ matter, since $L^{H'}\supseteq
L^H$ for $H'\le H$.  Writing $B_d$ for the matrix whose columns are the coordinate vectors of bases of
these fields, $a\in W_d(L)$ if and only if $B_dc=\operatorname{vec}(a)$ has a rational solution, and the
solution grouped by field gives the summands.

A failed membership test also has a small exact certificate: a rational row vector $\lambda$ with
$\lambda B_d=0$ but $\lambda\operatorname{vec}(a)\ne0$.  Such a separating functional exists precisely
when the system is inconsistent.  Checking it only requires rational matrix multiplication, although
the claim that it rules out every allowed summand still requires the field list to be complete.

\subsection{The permutation-module reformulation}

The criterion can be pushed down from $L$ to the $\Q$-span of the conjugates of $a$, which has dimension
at most $n$ and needs no basis of $L$.  Let $a_1,\dots,a_n$ be the conjugates of $a$ and
$V=\operatorname{span}_\Q\{a_1,\dots,a_n\}\subseteq L$, a $G$-stable subspace.  For $S\subseteq
\{1,\dots,n\}$ write $s_S=\sum_{i\in S}a_i$.

\begin{theorem}[conjugate-subset criterion, \cite{r1}]\label{thm:subsets}
$\Dp(a)\le d$ if and only if $a$ lies in the $\Q$-span of the subset sums $s_S$ with $\dg(s_S)\le d$.
\end{theorem}

\begin{proof}
Since $\Q[G]$ is semisimple, $V$ has a $G$-stable complement in $L$ and there is a $G$-equivariant
projection $P:L\to V$ (average any projection over $G$).  $P$ is the identity on $V$, so $P(a)=a$, and
$P(L^H)\subseteq V^H$ for every $H$.  Hence $a\in W_d(L)$ implies $a\in\sum_{[G:H]\le d}V^H$; the
converse holds because $V^H\subseteq L^H$.  Finally $V^H$ is spanned by the $H$-orbit sums of the $a_i$:
the surjection $\Q^n\to V$, $e_i\mapsto a_i$, is $G$-equivariant, so by averaging lifts over $H$ it
restricts to a surjection $(\Q^n)^H\to V^H$, and $(\Q^n)^H$ is spanned by the indicator vectors of the
$H$-orbits.  Each orbit sum is fixed by $H$ and so has degree at most $[G:H]\le d$.
\end{proof}

Thus the additive problem lives entirely in the permutation module of $G$ on the roots.  The only
information beyond the group that it needs is the space of $\Q$-linear relations among the conjugates
(the kernel of $\Q^n\to V$), which the implementation obtains from the trace form
(Section~\ref{sec:impl}).

\subsection{Examples}

\begin{example}[the sum example]
For $\alpha_+$ the splitting field is $L=L_fL_g$ of degree $36$ with $G=S_3\times S_3$.  The subgroup
counts by index are $1,3,6,1,20,9,4,15,1$ for indices $1,2,3,4,6,9,12,18,36$ (sixty subgroups).  The
six cubic subfields are the fields generated by the conjugates of $u$ and of $v$, and
$\alpha_+\in\Q(u)+\Q(v)$ recovers \eqref{eq:answers}.  Since $3$ is also the prime lower bound, the
answer is globally optimal.
\end{example}

\begin{example}[the product root as a sum]\label{ex:sextic}
The same field serves $\alpha_\times=uv$.  Inside $K=\Q(uv)$, whose subfields have degrees $1,3,3,9$
(they are $\Q$, $\Q(u)$, $\Q(v)$, $K$), the sum $\Q(u)+\Q(v)=\operatorname{span}\{1,u,u^2,v,v^2\}$ does
not contain $uv$ because $\{u^iv^j\}_{0\le i,j\le2}$ is a basis of $K$; so $\Dp^K(uv)=9$.  Globally,
label the roots $u_1=u,u_2,u_3$ and $v_1=v,v_2,v_3$ and put $w_\pi=\sum_ju_jv_{\pi(j)}$ for
$\pi\in S_3$.  The two permutations fixing $1$ give real numbers $w_1,w_2$ with
\[
 w_1+w_2=2uv+(u_2+u_3)(v_2+v_3)=3uv
\]
because both cubics have zero trace.  The six values $w_\pi$ are the roots of the irreducible sextic
\begin{equation}\label{eq:R}
 R(x)=x^6+6x^4-27x^3+9x^2-81x+4=\Bigl(x^3+3x-\tfrac{27}2\Bigr)^2-\tfrac{713}4,
\end{equation}
which has exactly two real roots; indeed $w_1,w_2$ satisfy $w^2-3uv\,w+3u^2-3v^2+4=0$ and eliminating
$u,v$ with the two cubics gives $R(w)^3$.  Hence
\begin{align*}
 \alpha_\times&=\tfrac13\Root[R,1]+\tfrac13\Root[R,2]\\
 &=\Root[S,1]+\Root[S,2],\qquad S(x)=R(3x),
\end{align*}
a sum of two sextics, and $\Dp(uv)\le6$.  To see that $6$ is optimal, let $A=C_3\times C_3\le G$ be
the Sylow $3$-subgroup, $U=\operatorname{span}\{u_1,u_2,u_3\}$ and $V=\operatorname{span}\{v_1,v_2,v_3\}$;
these are the two-dimensional standard representations of the two factors, and linear disjointness of
$L_f,L_g$ makes multiplication $U\otimes V\to L$ injective, so $T=\operatorname{span}\{u_iv_j\}$ is a
four-dimensional irreducible $G$-module containing $uv$.  A subgroup $H$ of index at most $5$ has index
$1,2,3$ or $4$.  If the index is $1,2$ or $4$ then $H\supseteq A$ (the index-$4$ subgroup is $A$
itself), and $A$ has no nonzero fixed vector in $T$.  If the index is $3$, then $|H|=12$, $H\cap A$ has
order $3$, and $H$ maps onto $G/A=C_2\times C_2$; the two independent sign changes act on
$A\cong\F_3^2$ by inverting one coordinate each, and a line of $\F_3^2$ invariant under both is a
coordinate axis, so $H\supseteq C_3\times1$ or $1\times C_3$, again giving $T^H=0$.  The equivariant
projection $L\to T$ kills every $L^H$ with $[G:H]\le5$ and fixes $uv$, so $uv\notin W_5(L)$.  Therefore
$\Dp(\alpha_\times)=6$, while $\Dp^{K}(\alpha_\times)=9$ and $\Dt(\alpha_\times)=3$.
\end{example}

\begin{example}[a sextic needing quartic summands outside its own field]\label{ex:quartic}
Let $r_1<r_2$ be the real roots of $x^4-x-1$ and $s=r_1+r_2$.  Then
$\Res_y(y^4-y-1,(x-y)^4-(x-y)-1)=(x^4-8x-16)(x^6+4x^2-1)^2$ and $s$ is the positive root of the
irreducible sextic $x^6+4x^2-1$.  The quartic is irreducible modulo $2$ and has pattern $(1,3)$ modulo
$7$, so its Galois group contains a $4$-cycle and a $3$-cycle and is $S_4$; the stabilizer of the pair
$\{r_1,r_2\}$ is $\langle(12),(34)\rangle$ of index $6$, whose only proper overgroup is the dihedral
group of order $8$.  So $\Q(s)$ has subfields of degrees $1,3,6$ only and $\Dp^{\Q(s)}(s)=6$; but
$s=r_1+r_2$ shows $\Dp(s)\le4$, and Theorem~\ref{thm:bound}(ii) with $\exp(S_4)=12\nmid\lcm(1,2,3)$
gives $\Dp(s)=4$.  For the stated real branches the identity is
\[
 \Root[-1+4\#^2+\#^6\&,2]
 =\Root[-1-\#+\#^4\&,1]+\Root[-1-\#+\#^4\&,2].
\]
\end{example}

\begin{remark}[real summands]\label{rem:real}
If $a$ is real and one insists on real summands, the criterion is obtained by restricting
\eqref{eq:Wd} to subgroups containing complex conjugation: trace projection commutes with conjugation,
so it preserves reality.
\end{remark}

\section{Binary splittings, flat sums, and Gaussian rational coefficients}\label{sec:binary}

A natural strategy for producing decompositions is to split $a$ into \emph{two} \Root{}s of smaller
degree and to iterate on the pieces.  The following question was raised while this article was
written: can a \Root{} be a linear combination of \Root{}s of smaller degree with Gaussian rational
coefficients, while no such iterated binary splitting exists?  The answer is yes, and the phenomenon
already occurs with rational coefficients.

\begin{definition}
A \emph{binary additive step} for $a$ is a representation $a=b+c$ with $\dg b,\dg c<\dg a$; a
\emph{binary multiplicative step} is $a=bc$ with the same condition.  A \emph{binary tree} for $a$
applies such steps recursively.  Its leaves are components of degree smaller than $\dg a$; the first
step alone already gives a decomposition into two components of smaller degree.
\end{definition}

\begin{proposition}[binary criteria]\label{prop:binary}
Let $L$ be the splitting field of $a$, $n=\dg a$.  A binary additive step exists if and only if
$a\in E_1+E_2$ for two subfields $E_1,E_2\subseteq L$ of degree less than $n$.  A binary multiplicative
step exists if and only if the two-factor criterion of Theorem~\ref{thm:two} holds with $d=n-1$.
\end{proposition}

\begin{proof}
For sums, apply Lemma~\ref{lem:trace} to $a=b+c$: the projections $\pi(b),\pi(c)$ lie in $L$, have
degrees at most $\dg b,\dg c<n$, and add up to $a$.  The converse is clear.  The product statement is
Theorem~\ref{thm:two}.
\end{proof}

\begin{theorem}\label{thm:flat}
Let $\eta=\sqrt2+\sqrt3+\sqrt6$, of degree $4$ with minimal polynomial $x^4-22x^2-48x-23$.  Then
$\Dp(\eta)=2$, attained by the flat sum $\sqrt2+\sqrt3+\sqrt6$, but $\eta$ admits no binary additive
step and no binary multiplicative step.  The same holds if Gaussian rational coefficients are allowed in
the binary steps, that is, there are no $c,c'\in\Q(i)$ and algebraic $X,Y$ of degree at most $3$ with
$\eta=cX+c'Y$ or $\eta=cXY$.
\end{theorem}

\begin{proof}
$\Q(\eta)=L=\Q(\sqrt2,\sqrt3)$ is Galois with group $C_2\times C_2$; its proper subfields are $\Q$
and the three quadratic fields $\Q(\sqrt2),\Q(\sqrt3),\Q(\sqrt6)$.  In the basis $1,\sqrt2,\sqrt3,
\sqrt6$ the vector-space sum of any two of these fields misses at least one of the three coordinates in
which $\eta$ is nonzero, so $\eta\notin E_1+E_2$ and there is no binary additive step by
Proposition~\ref{prop:binary}.  For products apply Theorem~\ref{thm:two} with $d=3$: the exponent
satisfies $t\le\min(3,\lfloor9/4\rfloor)=2$, and $t=2$ forces $E=F=\Q$, i.e.\ $\eta^2\in\Q$, which is
false.  For $t=1$ the fields $E,F$ are quadratic (or $\Q$), and one checks the three pairs directly: a
product $(x_0+x_1\sqrt2)(y_0+y_1\sqrt3)$ has coordinate vector $(x_0y_0,x_1y_0,x_0y_1,x_1y_1)$, whose
first coordinate vanishes only if the second or third does; the pairs $(\Q(\sqrt2),\Q(\sqrt6))$ and
$(\Q(\sqrt3),\Q(\sqrt6))$ are the same computation.  So $\Dtt(\eta)=4$.

Now allow Gaussian coefficients.  For sums, suppose $\eta=cX+c'Y$ with $c,c'\in\Q(i)$ and
$\dg X,\dg Y\le3$.  Work in the Galois field $L'=L(i)=\Q(i,\sqrt2,\sqrt3)$, of degree $8$.  The
projection $\pi:M\to L'$ of Lemma~\ref{lem:trace} is $L'$-linear, so $\eta=c\pi(X)+c'\pi(Y)$ with
$\pi(X),\pi(Y)\in L'$ of degree at most $3$, hence lying in $\Q$ or in a quadratic subfield $\Q(\sqrt m)$,
$m\in\{-1,\pm2,\pm3,\pm6\}$.  Take real parts: if $m>0$ then $\operatorname{Re}(cX')=\operatorname{Re}(c)
X'\in\Q(\sqrt m)$, and if $m<0$, $X'=p+qi\sqrt{|m|}$ with $p,q\in\Q$ and $c=e+e'i$ give
$\operatorname{Re}(cX')=ep-e'q\sqrt{|m|}\in\Q(\sqrt{|m|})$.  So $\eta$ is a sum of two elements of
real quadratic subfields of $L$, which was excluded above.  For products, $\eta=cXY$ with $c\notin\Q$
means that $\eta/c$ is a product of two numbers of degree at most $3$; writing $c^{-1}=e+e'i$ with
$e'\ne0$, the element $\eta/c=e\eta+e'i\eta$ has nonzero coordinates on $i\sqrt2,i\sqrt3,i\sqrt6$.
An automorphism sending $(i,\sqrt2,\sqrt3)$ to $(\epsilon_0i,\epsilon_1\sqrt2,\epsilon_2\sqrt3)$
can fix it only if $\epsilon_0\epsilon_1=\epsilon_0\epsilon_2=
\epsilon_0\epsilon_1\epsilon_2=1$, which forces every $\epsilon_j=1$.  Thus its stabilizer in
$\Gal(L'/\Q)$ is trivial and it has degree $8$.  Theorem~\ref{thm:two} with $n=8$,
$d=3$ allows only $t=1$ and quadratic $E,F$, whose compositum has degree at most $4<8$.  If $c\in\Q$
the rational case applies.
\end{proof}

The same argument shows that $i+\sqrt2+i\sqrt2$, of degree $4$, is the two-term Gaussian combination
$i\cdot1+(1+i)\sqrt2$ of \Root{}s of degree $1$ and $2$ but admits no binary step with rational
coefficients.  The general principle is that a binary step tests membership in the sum of \emph{two}
subfields, whereas Theorem~\ref{thm:sum} allows the sum of all of them; a flat sum of $k$ terms can
require $k$ distinct subfields none of whose pairwise sums contains $a$.

\begin{theorem}[Gaussian coefficients, flat sums]\label{thm:gaussian}
Let $L'$ be a finite Galois extension containing $a$ and $i$, with group $G'$.  Then $a=\sum_jc_jb_j$
with $c_j\in\Q(i)$ and $\dg(b_j)\le d$ for some finite $r$ if and only if
\[
 a\in\sum_{H\le G',\,[G':H]\le d}\bigl(L'^H+i\,L'^H\bigr).
\]
\end{theorem}

\begin{proof}
The projection $\pi:M\to L'$ is $L'$-linear, hence $\Q(i)$-linear, so $a=\sum c_j\pi(b_j)$ with
$\pi(b_j)$ in subfields $E_j$ of $L'$ of degree at most $d$, and $c_jE_j\subseteq E_j+iE_j$.  Conversely
$E+iE$ consists of the two-term combinations $1\cdot e+i\cdot e'$.
\end{proof}

The Gaussian flat problem is therefore again exact linear algebra (the package option
\code{"Coefficients" -> "GaussianRationals"} adjoins $i$ to the splitting field).  The lower bounds
change: adjoining Gaussian coefficients adds a factor $C_2$ to the product of normal-closure groups
in Theorem~\ref{thm:bound}.  Consequently only prime divisors $\ge3$ of $n$ give the prime bound, while
the sharper exponent condition is
\[
 \exp(G')\mid\lcm\bigl(2,\lcm(1,\dots,d)\bigr),
\]
when $L'$ is the splitting field of $a$ with $i$ adjoined.  Indeed this field is contained in the
compositum of $\Q(i)$ with the normal closures of the components.  For $d\ge2$ the exponent bound
is unchanged; Gaussian coefficients only waive its obstruction at $d=1$.  The binary Gaussian problem $a=cX+c'Y$ is, by
contrast, not linear ($cE$ is a subspace for each fixed $c$, but their union over $c$ is not); the
real-part argument above settles it for real targets in totally real multiquadratic fields, and the
package answers the rational binary question through the option \code{"MaxTerms" -> 2}.

\section{Products of two factors: a complete algorithm}\label{sec:products}

Trace is additive; norm is multiplicative but takes $a$ to a power $a^t$.  Keeping that power is what
makes the following criterion complete.

\begin{lemma}[intersection lemma]\label{lem:intersection}
Let $L/\Q$ be finite Galois and $b$ algebraic.  Put $E_0=L\cap\Q(b)$ and $t=[L(b):L]$.  Then
$\dg(b)=t\,[E_0:\Q]$, the minimal polynomial of $b$ over $L$ has coefficients in $E_0$, and
$\Nm_{L(b)/L}(b)\in E_0$.
\end{lemma}

\begin{proof}
Let $M/\Q$ be Galois containing $L$ and $b$, $\Gamma=\Gal(M/\Q)$, $N=\Gal(M/L)\trianglelefteq\Gamma$,
$J=\Gal(M/\Q(b))$.  The compositum $L(b)$ corresponds to $N\cap J$ and the intersection $E_0$ to $NJ$,
which is a subgroup because $N$ is normal.  Hence $[\Q(b):E_0]=[NJ:J]=[N:N\cap J]=[L(b):L]=t$ and
$\dg(b)=t[E_0:\Q]$.  The roots of the minimal polynomial of $b$ over $L$ form the $N$-orbit of $b$;
an element of $J$ fixes $b$ and permutes this orbit (normality of $N$), so the coefficients of that
polynomial, being symmetric functions of the orbit, are fixed by $NJ$ and lie in $E_0$.  The norm is
the constant coefficient up to sign.
\end{proof}

\begin{theorem}[two-factor criterion]\label{thm:two}
Let $L/\Q$ be finite Galois, $0\ne a\in L$, $n=\dg(a)$, and $d\ge1$.  There exist algebraic $b,c$ with
$a=bc$ and $\dg b,\dg c\le d$ if and only if there are subfields $E,F\subseteq L$ and an integer $t\ge1$
with
\begin{equation}\label{eq:two}
 t[E:\Q]\le d,\qquad t[F:\Q]\le d,\qquad E\cap a^tF\ne\{0\}.
\end{equation}
It suffices to consider $1\le t\le\min(d,\lfloor d^2/n\rfloor)$, and pairs with
$n\le t[E:\Q][F:\Q]$.
\end{theorem}

\begin{proof}
Necessity.  If $a=bc$ then $L(b)=L(c)$ because $c=a/b$ and $b=a/c$; call it $M$ and put $t=[M:L]$,
$E=L\cap\Q(b)$, $F=L\cap\Q(c)$.  By Lemma~\ref{lem:intersection}, $\dg b=t[E:\Q]\le d$ and
$\dg c=t[F:\Q]\le d$, and $u=\Nm_{M/L}(b)\in E$, $v=\Nm_{M/L}(c)\in F$ are nonzero with
$uv=\Nm_{M/L}(a)=a^t$.  So $u=a^t/v\in E\cap a^tF$.

Sufficiency.  Let $0\ne u\in E\cap a^tF$, say $u=a^tw$ with $w\in F$.  Choose any $t$th root $b$ of
$u$ and put $c=a/b$.  Then $b^t=u\in E$ and $c^t=a^t/u=1/w\in F$, so $b$ and $c$ are roots of
$X^t-u\in E[X]$ and $X^t-1/w\in F[X]$ and $\dg b\le t[E:\Q]\le d$, $\dg c\le t[F:\Q]\le d$.  The
product is $a$ by construction; no root-of-unity ambiguity arises because $c$ is defined by division.

Bounds.  If \eqref{eq:two} holds then $a^t\in EF$, so $n\le t\dg(a^t)\le t[EF:\Q]\le t[E:\Q][F:\Q]
\le d^2/t$; and $t\le d$ since $t[E:\Q]\le d$.
\end{proof}

The sharper necessary inequality $n\le t[EF:\Q]$ is useful computationally.  If $E=L^H$ and $F=L^J$,
then $EF=L^{H\cap J}$ and $[EF:\Q]=[G:H\cap J]$, so intersecting the two finite subgroups can reject
a pair before any rational null-space calculation.  This improves the product-of-degrees bound when
the fields overlap, including all nested pairs.

\begin{corollary}\label{cor:two}
$\Dtt(a)$ is computable: test $d=\max(\lceil\sqrt n\rceil,\text{lower bounds}),\dots,n-1$ in turn, for
each $d$ the exponents $t$ and all pairs of subfields of $L$; the first success gives $\Dtt(a)$ and a
decomposition, and if none succeeds then $\Dtt(a)=n$.  With $t=1$ and both fields inside a fixed number
field $K$ (not necessarily Galois), the same test decides whether $a\in E^\times F^\times$, hence
computes the two-factor minimum inside $K$.  If $K$ is Galois, testing \emph{all} the permitted
exponents computes $\Dtt(a)$ globally; restricting to $t=1$ still only decides the problem inside $K$.
\end{corollary}

The condition $E\cap a^tF\ne\{0\}$ is rational linear algebra: with column bases $B_E,B_F$ and the
matrix $M_{a^t}$ of multiplication by $a^t$, it holds if and only if $[\,B_E\ \ -M_{a^t}B_F\,]$ has a
nonzero null vector $(x,y)$; then $u=B_Ex\ne0$ (else $y=0$ too), and $b=u^{1/t}$, $c=a/b$.  No unit
groups, ideal factorizations or logarithms are needed for two factors.  For the product example
\eqref{eq:examples} the pair $(\Q(u),\Q(v))$ with $t=1$ succeeds at $d=3$, which is also the lower
bound; hence $\Dt(\alpha_\times)=\Dtt(\alpha_\times)=3$.

\subsection{Optimal factors outside the splitting field}\label{sec:external}

\begin{theorem}\label{thm:external}
Let $a=\sqrt{(1+\sqrt2)(1+\sqrt3)}>0$, with minimal polynomial $x^8-4x^6-16x^4-8x^2+4$
(it is \code{Root[\ldots,4]}).  Then $\Dt(a)=\Dtt(a)=4$, attained by
$a=\sqrt{1+\sqrt2}\cdot\sqrt{1+\sqrt3}$ with quartic factors of minimal polynomials $x^4-2x^2-1$ and
$x^4-2x^2-2$, and $\Dp(a)=8$.  Moreover no finite sum or product of elements of the splitting field
$L$ of $a$ of degree less than $8$ equals $a$; the optimal factors lie outside $L$, and at the optimal
bound $d=4$ the criterion \eqref{eq:two} holds only with $t=2$.
\end{theorem}

\begin{proof}
Write $r=\sqrt2$, $s=\sqrt3$, $A=1+r$, $B=1+s$, $B_0=\Q(r,s)$.  The four conjugates of $a^2=AB$ are
distinct, so $\Q(a^2)=B_0$; $AB$ is negative at some real embedding of the totally real field $B_0$,
so it is not a square there and $\dg a=8$.  Put $B_1=B_0(i)$ and $L=B_1(a)$.  Since $a$ is real and
$B_1\cap\R=B_0\not\ni a$, $[L:\Q]=16$.  The conjugates of $a$ are $\pm a,\pm iB/a,\pm irA/a,\pm r/a$: each squares to a conjugate of
$AB$, for instance $(iB/a)^2=-B^2/(AB)=-B/A=B(1-r)=(1+s)(1-r)$ because $(1+r)(1-r)=-1$; all of them
lie in $L$; conversely the splitting field contains $a\cdot iB/a=iB$, hence $i$, and $B_0$, so
$L$ is exactly the splitting field.  Let $z$ fix $B_1$ and send $a\mapsto-a$; let $\sigma,\tau,\kappa$
be the involutions acting by $(r,s,i,a)\mapsto(-r,s,i,iB/a)$, $(r,-s,i,irA/a)$ and $(r,s,-i,a)$ (complex
conjugation).  Direct substitution shows $z$ is central, $[\sigma,\tau]=[\sigma,\kappa]=[\tau,\kappa]=z$,
every element is uniquely $\sigma^e\tau^f\kappa^gz^h$ with exponents in $\{0,1\}$, and
$(\sigma^e\tau^f\kappa^gz^h)^2=z^{ef+eg+fg}$; so $G=\Gal(L/\Q)$ has order $16$ and exponent $4$.

\emph{Every subgroup of order at least $4$ contains $z$.}  Elements of order $4$ square to $z$.  A
subgroup of order $4$ avoiding $z$ would be a Klein four-group of non-central involutions; these lie in
the three cosets $\sigma\langle z\rangle,\tau\langle z\rangle,\kappa\langle z\rangle$, involutions from
different cosets do not commute, and two from the same coset multiply to $z$.  Every larger $2$-group
contains a subgroup of order $4$.

Consequently every subfield of $L$ of degree $1,2$ or $4$, i.e.\ of degree less than $8$, is contained
in $L^{\langle z\rangle}=B_1$, and since $a\notin B_1$ no sum or product of such elements is $a$.  Trace
projection (Lemma~\ref{lem:trace}) moves any additive decomposition into $L$ without increasing
degrees, so $\Dp(a)=8$.  The displayed factorization gives $\Dt(a)\le4$, and Theorem~\ref{thm:bound}(ii)
with exponent $4\nmid\lcm(1,2,3)=6$ gives $\Dt(a)\ge4$.  In the criterion, $E=\Q(\sqrt2)$,
$F=\Q(\sqrt3)$, $t=2$, $u=1+\sqrt2$ works: $a^2/u=1+\sqrt3\in F$.  Every pair with $t=1$ fails because
the fields of degree at most $4$ lie in $B_1$.
\end{proof}

The reports \cite{r1,r3,r7} point out that Theorem~3.1 of the manuscript \cite{Altamirano}, as they
read it, asserts that two lower-degree factors can always be replaced by lower-degree factors inside
the normal closure, and that its proof divides by a coefficient that can vanish (the linear coefficient
of $X^2-(1+\sqrt2)$ here).  We have not verified the manuscript independently; Theorem~\ref{thm:external}
shows that any such descent statement without the exponent $t$ is false.  A related norm-and-radical
replacement, with an explicit root of unity, appears in Samuels's work on the metric Mahler measure
\cite{Samuels}; Mahler measure ignores roots of unity, absolute degree does not, so that result does not
transfer.

\section{Products of arbitrarily many factors}\label{sec:many}

\subsection{Two factors are not enough}

\begin{example}\label{ex:eta}
$\eta_3=(1+\sqrt2)(1+\sqrt3)(1+\sqrt5)$ has degree $8$ (all eight coordinates in the multiquadratic
basis are nonzero, so no sign change fixes it) and $\Dt(\eta_3)=2$ by the prime bound.  But two factors
of degree at most $3$ cannot work: the compositum of two fields of degree at most $3$ has degree in
$\{1,2,3,4,6,9\}$, never a multiple of $8$, so it cannot contain $\eta_3$.  Hence $\Dtt(\eta_3)=4$,
attained e.g.\ by $(1+\sqrt2)(1+\sqrt3)\cdot(1+\sqrt5)$.
\end{example}

\begin{example}\label{ex:w}
$w=(1+\sqrt2)(1+\sqrt3)(1+\sqrt6)=7+4\sqrt2+3\sqrt3+2\sqrt6$ has degree $4$ and minimal polynomial
$x^4-28x^3+128x^2-200x+100$.  $\Dt(w)=2$, but $\Dtt(w)=4$: in $L=\Q(\sqrt2,\sqrt3)$ with $d=3$,
the only permitted exponents are $t=1,2$.  The case $t=2$ forces $E=F=\Q$ and $w^2\in\Q$, which is
false.  For $t=1$, the three pairs of quadratic subfields are excluded because the
product $(x_0+x_1\sqrt m)(y_0+y_1\sqrt{m'})$ has rank-one coefficient matrix, while the matrices of
$w$ are
$\begin{psmallmatrix}7&3\\4&2\end{psmallmatrix}$,
$\begin{psmallmatrix}7&2\\4&3/2\end{psmallmatrix}$,
$\begin{psmallmatrix}7&2\\3&4/3\end{psmallmatrix}$
with determinants $2,\tfrac52,\tfrac{10}3$.  Here even a greedy strategy fails: no factorization of
$w$ into two factors of smaller degree exists at all, yet three quadratic factors do.
\end{example}

\begin{example}
For $c$ the real root of $x^3+x+3$, the product $uvc$ of three cubic roots has degree $27$ and
$\Dt=3$ (prime bound); no two factors of degree $\le3$ can reach degree $27>9$ \cite{r2}.
\end{example}

\subsection{Norms, radicals and the root-of-unity obstruction}

\begin{proposition}[necessary condition]\label{prop:normnec}
Let $L/\Q$ be Galois, $0\ne a\in L$, and $a=b_1\cdots b_r$ with $\dg b_i\le d$.  Put $E_i=L\cap\Q(b_i)$,
$t_i=[L(b_i):L]$, $n_i=\Nm_{L(b_i)/L}(b_i)\in E_i$ and $T=[L(b_1,\dots,b_r):L]$.  Then $t_i[E_i:\Q]\le d$,
$t_i\mid T$, and
\begin{equation}\label{eq:normnec}
 a^T=\prod_{i=1}^rn_i^{T/t_i}.
\end{equation}
\end{proposition}

\begin{proof}
Take norms from $M=L(b_1,\dots,b_r)$ to $L$: $\Nm_{M/L}(b_i)=\Nm_{L(b_i)/L}(b_i)^{[M:L(b_i)]}$ by
transitivity, and $\Nm_{M/L}(a)=a^T$.  The degree statements are Lemma~\ref{lem:intersection}.
\end{proof}

\begin{proposition}[sufficient condition]\label{prop:radsuff}
If there are subfields $E_i\subseteq L$, integers $t_i\ge1$ with $t_i[E_i:\Q]\le d$, and $n_i\in E_i^\times$
with $a=\prod_in_i^{1/t_i}$ for some choice of the radicals, then $\Dt(a)\le d$.
\end{proposition}

\begin{proof}
$b_i=n_i^{1/t_i}$ satisfies $X^{t_i}-n_i\in E_i[X]$, so $\dg b_i\le t_i[E_i:\Q]\le d$.
\end{proof}

The gap between the two propositions is a root of unity.  From \eqref{eq:normnec}, the radicals
$b_i'=n_i^{1/t_i}$ satisfy $\prod b_i'=\zeta a$ with $\zeta^T=1$; if $\tau=\lcm(t_i)$ then $\zeta^\tau=
\prod n_i^{\tau/t_i}/a^\tau\in L$, so $\zeta$ is a root of unity of order dividing $\tau\,|\mu(L)|$.
The radicals may be adjusted by $t_i$th roots of unity, which absorb $\zeta$ when $\zeta\in\mu_\tau$,
but not in general: $\zeta_8=(1+i)/\sqrt2$ is a product of numbers of degree $2$ although $\dg\zeta_8=4$,
so a stray root of unity cannot simply be appended as a low-degree factor.  For two factors this issue
does not arise ($t_1=t_2=T$, and $c$ is defined by division), which is why Theorem~\ref{thm:two} is
complete.

\begin{problem}
Construct a terminating decision procedure for $\Dt(a)$.  Equivalently, given a Galois number field $L$, $a\in L^\times$ and $d$, decide
whether $a$ lies in the subgroup of $\Qbar^\times$ generated by the multiplicative groups of all number
fields of degree at most $d$.  Even the restricted question with all factors in $L$ and $t=1$ --- is
$a\in\prod_{[E:\Q]\le d}E^\times$? --- is a question about multiplicative groups of number fields
(units, class groups, $S$-units), not about vector spaces.  No complete criterion for either
many-factor problem is established here, and this formulation is not a claim that their decidability
is an established open problem in the literature.  A positive semi-decision procedure exists by
enumerating the number of factors, degrees, coefficient heights and branches and checking the resulting
identities exactly; the present methods do not supply a terminating negative certificate in general.
\end{problem}

\subsection{Arithmetic obstructions beyond Galois theory}

\begin{theorem}[sign obstruction, \cite{r4}]\label{thm:sign}
Let $\theta=1+\sqrt2+\sqrt3$, of degree $4$ with minimal polynomial $x^4-4x^3-4x^2+16x-8$ and norm
$\Nm_{\Q(\sqrt2,\sqrt3)/\Q}(\theta)=-8$.  Then $\Dp(\theta)=2$ and $\Dt(\theta)=4$.
\end{theorem}

\begin{proof}
$\Dp(\theta)\le2$ is visible and $\theta\notin\Q$.  For products we show that no finite product of
numbers of degree at most $3$ equals $\theta$.

\emph{Step 1: cubic factors can be removed after an odd power.}  Let $\theta=\prod\beta_j$ with
$\dg\beta_j\le3$.  Let $K$ be the compositum of $L=\Q(\sqrt2,\sqrt3)$, the quadratic fields of the
quadratic $\beta_j$, and the quadratic resolvent fields of the cubic $\beta_j$; $K$ is multiquadratic,
of degree a power of two.  A cubic $\beta_j$ has $[K(\beta_j):K]=3$ (a cubic with a root in a field of
$2$-power degree is impossible), its minimal polynomial over $K$ is its rational cubic, and
$\Nm_{K(\beta_j)/K}(\beta_j)=\Nm_{\Q(\beta_j)/\Q}(\beta_j)\in\Q^\times$.  Let $E$ be the compositum of
$K$ with the splitting fields of the cubics; each is cyclic of degree $3$ over $K$ because the
resolvent is already in $K$, so $[E:K]=3^s=:m$ is odd.  Taking $\Nm_{E/K}$ of the product gives
$\theta^m=q\prod_{\dg\beta_j\le2}\beta_j^m$ with $q\in\Q^\times$: an odd power of $\theta$ is a product
of numbers of degree at most $2$.

\emph{Step 2: a product of quadratics in a totally real biquadratic field has positive norm.}  Let
$\theta'\in L$ be a finite product of numbers of degree $\le2$.  The non-real quadratic factors have a
real product (the other factors and $\theta'$ are real), and each non-real quadratic $\beta$ has
$\beta\bar\beta\in\Q_{>0}$, so the square of their product is a positive rational and they may be
replaced by a single real quadratic factor.  Now all factors are real; let $M\supseteq L$ be a totally
real multiquadratic field containing them.  For a real quadratic $v$ the function
$\sigma\mapsto\operatorname{sign}(\sigma v)$ on $\Gal(M/\Q)$ is a constant times a character with values
$\pm1$, and this form is preserved under products.  Since $\theta'\in L$, the resulting function is
constant on cosets of $\Gal(M/L)$ and descends to $\Gal(L/\Q)=C_2\times C_2$, where a constant has
$0$ or $4$ negative values and a nontrivial character exactly $2$; so the product of the signs of the
four conjugates of $\theta'$ is $+1$, i.e.\ $\Nm_{L/\Q}(\theta')>0$.

Applying Step 2 to $\theta'=\theta^m$ gives $\Nm(\theta^m)=(-8)^m<0$, a contradiction.  Hence no
product of numbers of degree $\le3$ equals $\theta$, and $\Dt(\theta)=4$.
\end{proof}

\begin{corollary}\label{cor:negative-norm}
If $L$ is a totally real biquadratic field and $a\in L$ has $\Nm_{L/\Q}(a)<0$, then $\Dt(a)=4$.
\end{corollary}

\begin{proof}
The same odd-power descent and sign argument excludes every product of numbers of degree at most
$3$.  The one-factor representation gives the opposite bound, since $\dg(a)\le[L:\Q]=4$.
\end{proof}

Here the Galois group is $C_2\times C_2$: the prime bound and the exponent bound give $2$, yet
$\Dt=4$.  Any complete multiplicative algorithm must see arithmetic (here the sign of a norm), which is
why linear algebra alone cannot solve the general many-factor problem.

\subsection{An exact criterion for linearly disjoint families}

\begin{proposition}[rank-one tensor test, \cite{r9}]\label{prop:tensor}
Let $K$ be a number field, $a\in K^\times$, and $E_1,\dots,E_r\subseteq K$ subfields with $\prod_i[E_i:\Q]=[K:\Q]$ such
that the products of basis elements span $K$ (equivalently, the multiplication map
$E_1\otimes_\Q\cdots\otimes_\Q E_r\to K$ is an isomorphism).  Write $a=\sum c_{j_1\cdots j_r}
e_{1,j_1}\cdots e_{r,j_r}$ in the product basis.  Then $a=u_1\cdots u_r$ with $u_i\in E_i$ if and only
if the tensor $c$ has rank one, which is decided by the pivot test: for a nonzero entry $c_p$, $c$ has
rank one if and only if
\[
 c_{j_1\cdots j_r}=c_p\prod_{i=1}^r\frac{c_{p(i\leftarrow j_i)}}{c_p}
 \qquad\text{for every }(j_1,\dots,j_r),
\]
where $p(i\leftarrow j)$ is $p$ with its $i$th index replaced by $j$.  The factors are then read off
from the slices through $p$, absorbing the leading scalar $c_p$ into one factor.  For $a=0$ the
decomposition is immediate and no nonzero pivot is needed.
\end{proposition}

\begin{proof}
The multiplication map is an isomorphism of $\Q$-spaces, so $a$ is a product of elements of the $E_i$
if and only if its tensor is a pure tensor, and the pivot formula characterizes pure tensors.
\end{proof}

For $\eta_3$ of Example~\ref{ex:eta}, $K=\Q(\sqrt2,\sqrt3,\sqrt5)$ and the three quadratic subfields
satisfy the hypothesis; the tensor is $(1,1)\otimes(1,1)\otimes(1,1)$ and the three factors are
recovered.  For $w$ of Example~\ref{ex:w} the three quadratic fields have $2\cdot2\cdot2\ne4$, the
hypothesis fails, and indeed the test does not apply; there a bounded dictionary search finds the
three factors.

\subsection{What the implementation does for many factors}

In order of cost: the complete two-factor algorithm; the tensor test over all families of subfields of
$L$ (or of $K$) whose degrees multiply to the ambient degree; recursive splitting of the factors
found (a heuristic: a two-factor optimum need not refine to a many-factor optimum, but it often does,
as for $\eta_3=[(1+\sqrt2)(1+\sqrt3)]\cdot(1+\sqrt5)$); and a small bounded dictionary search.  A
result is marked globally optimal only when its maximum degree equals a proved lower bound, and
\emph{two-factor optimal} when the complete two-factor search has been exhausted.

\section{Comparison of the nine reports}\label{sec:reports}

All nine reports establish the same core: the resultant identities, irreducibility and branch
selection for the examples, the largest-prime bound, the trace descent and the fixed-space criterion
for sums, the intersection criterion $E\cap aF\ne0$ for two factors in a fixed field, and the
statement that arbitrary-length products are not covered by a complete algorithm.  Table~\ref{tab:reports}
records what each report contributed beyond that.  Every report states that its Wolfram code was not
executed (the evaluator they had access to returned HTTP~404); the Python/SymPy checks they ran verify
the algebraic identities, not the packages.

\begingroup\small
\begin{longtable}{@{}p{30pt}p{\dimexpr\textwidth-42pt\relax}@{}}
\caption{What the nine reports added to the common core.}\label{tab:reports}\\
\toprule
Report & Distinctive contributions and remarks\\
\midrule
\endfirsthead
\toprule
Report & Distinctive contributions and remarks (continued)\\
\midrule
\endhead
\cite{r1} & Conjugate-subset criterion (Theorem~\ref{thm:subsets}) with the equivariant-projection
proof; a degree-eight counterexample $\sqrt{1+\sqrt2}\sqrt{1+\sqrt7}$ with group of order $32$ (an
extraspecial central product) and the same conclusions as Theorem~\ref{thm:external}; template
(Gr\"obner) recovery of the cubics.\\
\cite{r2} & Residual-resultant search with a height bound on one polynomial only; the complete
two-factor criterion with exponent $k$; the degree-27 example of three cubics; the $2^{1/4}$ example
($\Dp=4$).\\
\cite{r3} & Two-factor criterion with exponent $t$ and the pruning bound $n\le t[E:\Q][F:\Q]$; the
example $w=(1+\sqrt2)(1+\sqrt3)(1+\sqrt6)$ with $\Dtt=4>\Dt=2$ and the determinant proof; a
vanishing-linear-coefficient example $\sqrt{\sqrt2(1+\sqrt3)}=2^{1/4}\sqrt{1+\sqrt3}$; subfields by
coefficient fields of divisors.\\
\cite{r4} & Inverse composition with a fixed first polynomial; the sign obstruction
(Theorem~\ref{thm:sign}); the $x^4-x-1$ pair-sum example (Example~\ref{ex:quartic}); the $\zeta_5$
example; fixed fields by iterated intersection with $\ker(A_\sigma-I)$.\\
\cite{r5} & Principal subfields $V_g=\{h(\theta):g\mid h(x)-h(\theta)\}$ with a self-contained proof
(after \cite{SvH}); rational normalization of the recovered pair; explicit recovery formulas.\\
\cite{r6} & Subfields by coefficient fields of divisors; trace centering; the sextic additive
optimum $\Dp(uv)=6$ with the projection $T=(1-e_1)(1-e_2)$ argument; the multiplicative subgroup
formulation $B_d(K)$.\\
\cite{r7} & The degree-eight example of Theorem~\ref{thm:external} with a complete group-theoretic
proof; the exponent bound with $\lcm(1,\dots,d)$; the three-quadratic example $\eta_3$; the
many-factor norm relation \eqref{eq:normnec} and the $\zeta_8$ remark; the most complete package
design (never executed).\\
\cite{r8} & The nonmembership certificate (a separating functional $w^\top B_d=0$, $w^\top a\ne0$);
the $\sqrt{2+\sqrt2}$ example; the necessary norm-descent condition (saturation of $B_d(L)$).\\
\cite{r9} & Rank-one tensor test (Proposition~\ref{prop:tensor}); polynomial recovery certificates
$A_\times,B_\times,A_+$; the sextic pair-product example $r_1r_2$ with $\Dt^{\Q(r_1r_2)}=6>4=\Dt$; the
example $\sqrt2+\sqrt3+\sqrt6$ as a three-term sum not reducible to two terms.\\
\bottomrule
\end{longtable}
\endgroup

Points on which the reports are inconsistent with one another, or which the present work corrects:
\begin{itemize}
\item The reports describe the degree-eight example either as a counterexample to a published
  descent theorem \cite{r1,r7} or only as a gap in its proof \cite{r3}.  The mathematical content
  (Theorem~\ref{thm:external}) is not in dispute; we have phrased the attribution cautiously.
\item All reports construct the splitting field with \code{ToNumberField[roots, All]} and convert
  elements with \code{ToNumberField[a, theta]}.  On the examples this takes 14 minutes for the first
  call and 10 minutes per element conversion in Wolfram 15.0.1 (Section~\ref{sec:impl}); it is not a
  usable design.
\item Several reports enumerate subfields of a Galois field by subsets of the linear factors of the
  minimal polynomial ($2^{[L:\Q]-1}$ subsets) or by all subgroups without first computing the group;
  the numerical-resolvent approach below is far cheaper.
\item Report \cite{r9} notes that the binary product search of its package is complete only inside
  $\Q(a)$ and that $\sqrt2+\sqrt3+\sqrt6$ needs three terms; Section~\ref{sec:binary} strengthens
  this to a theorem covering both operations and Gaussian coefficients.
\end{itemize}

\section{Algorithms and implementation}\label{sec:impl}

\subsection{Why not the power basis of a primitive element}

All reports propose: compute all roots, call \code{ToNumberField[roots, All]} to obtain a primitive
element $\theta$ of $L$ with the roots as polynomials in $\theta$, factor the minimal polynomial of
$\theta$ over $L$ to get the automorphisms, then enumerate subgroups.  For the degree-nine examples
($[L:\Q]=36$) the first call took about 14 minutes in Wolfram 15.0.1 and a single
\code{ToNumberField[root, theta]} conversion a further 10 minutes.  The obstacle is not the degree
but the power basis: a generic primitive element $\theta=\sum w_i\alpha_i$ has a minimal polynomial with
coefficients around $10^{20}$, and the coordinates of the roots in the basis $1,\theta,\dots,\theta^{35}$
have enormous height.  By contrast, \code{RootReduce[theta + w*Root[p,k]]}, which adjoins one root to a
partial primitive element, takes about a second per step.  The implementation therefore never uses the
power basis of $\theta$.

\subsection{The Galois group by numerical resolvents}

Let $P$ be monic with integer coefficients (replace $a$ by $c\,a$, $c$ the leading coefficient of its
minimal polynomial, so that all roots are algebraic integers) and let $\beta_1,\dots,\beta_n$ be its
roots, computed to high precision.  Choose integer weights $w_1,w_2,\dots$ and consider the partial
sums $\theta_k=\sum_{i\le k}w_i\beta_i$.  The Galois group $G\le S_n$ is determined level by level: the
set $O_k$ of $k$-tuples $\tau$ with $(\beta_{\tau(1)},\dots,\beta_{\tau(k)})$ in the $G$-orbit of
$(\beta_1,\dots,\beta_k)$ is exactly the set of $\tau$ for which $\theta_\tau=\sum_{i\le k}
w_i\beta_{\tau(i)}$ is a root of the minimal polynomial $m_k$ of $\theta_k$, provided the map
$\tau\mapsto\theta_\tau$ is injective on the candidate set $\{(\tau',j):\tau'\in O_{k-1},\,j\notin\tau'\}$.
The candidate set is $G$-stable, so $F_k(x)=\prod_{\text{candidates}}(x-\theta_\tau)$ is a monic integer
polynomial; it is squarefree exactly when the injectivity holds, and then $m_k$ is its irreducible
factor vanishing at $\theta_k$.  At level $n$, $O_n=G$ as a set of permutations, and closure under
composition is verified.

The Python implementation does this rigorously in Arb ball arithmetic \cite{Arb}: $F_k$ is computed
as a ball polynomial, each coefficient ball must contain a unique integer (otherwise the precision is
doubled), $F_k$ is factored exactly over $\Z$ with FLINT \cite{FLINT}, and $\tau\in O_k$ is tested by
$m_k(\theta_\tau)\ni0$; since every true orbit element passes this test and the number of passing
candidates is checked to equal $\deg m_k$, the orbit is determined exactly.  The Wolfram implementation
obtains $m_k$ from \code{RootReduce} of $\theta_{k-1}+w_k\,\code{Root[P,k]}$ (exact), locates the roots
numerically and applies the same counting checks; a degenerate weight (degree drop or collision) is
detected and redrawn.  The dominant cost is the exact minimal polynomial (or the exact factorization
of $F_k$, of degree at most $n|G|$); the method is practical for $|G|$ up to a few hundred.

\subsection{The splitting field in a tower basis with trace coordinates}

The levels at which $\deg m_k$ grows give a tower $L=\Q(\beta_{g_1},\dots,\beta_{g_s})$ with relative
degrees $e_1,\dots,e_s$, $\prod e_j=|G|=N$, and the monomials $b=\prod_j\beta_{g_j}^{a_j}$,
$0\le a_j<e_j$, form a $\Q$-basis of $L$ consisting of algebraic integers of small height.  Every
$\sigma\in G$ acts on such a monomial by permuting the roots, so the values $\sigma(b)$ are known
numerically for all $\sigma$ and all basis elements: an $N\times N$ matrix $\mathrm{Val}$.  The trace
form $\Tr_{L/\Q}(b_jb_l)=\sum_\sigma\sigma(b_j)\sigma(b_l)$ is an integer matrix $\mathrm{Gram}=
\mathrm{Val}^\top\mathrm{Val}$ (with a transpose, not a conjugate transpose), nondegenerate since $L/\Q$ is separable, and the
coordinates of any $y\in L$ given by its conjugate vector $(\sigma(y))_\sigma$ are
$\mathrm{Gram}^{-1}(\Tr(yb_l))_l$.  When $y$ is integral these traces are integers; for a rational
combination of the integral basis elements, first clear its coordinate denominators, recover integer
traces, then divide back.  A unique integer in each Arb ball certifies the recovery in Python; ordinary
high-precision rounding in Wolfram has the qualification described below.  This yields, by rational
linear algebra with integer inputs: the coordinates of all roots, the matrix $A_\sigma$ of every
automorphism (conjugate vector of $\sigma(b_j)$ is $\tau\mapsto(\tau\sigma)(b_j)$, read from the
multiplication table of $G$), multiplication matrices, fixed fields $L^H=\bigcap_{h\in H}
\ker(A_h-I)$ for all subgroups (enumerated by closure from cyclic subgroups), and the linear systems of
Theorems~\ref{thm:sum} and~\ref{thm:two}.  An element with coordinate vector $v$ is converted back to
a \Root{} object by computing its distinct conjugates (the orbit of $v$ under the $A_\sigma$), forming
their polynomial after clearing coordinate denominators, and identifying the intended branch.  In
Python the branch must be the unique isolated root compatible with the element's enclosure;
precision is increased when the match is ambiguous.  A consistency check ($A_\sigma$ permutes the
root coordinates as $\sigma$ does) guards the numerical steps in Wolfram.  Its final identities and
degrees are checked by \code{RootReduce} and \code{MinimalPolynomial}; Python also supplies an
independent exact verifier using composed polynomials and certified branch isolation.

\subsection{The fast path inside \texorpdfstring{$\Q(a)$}{Q(a)}}

Independently of the Galois group, all subfields of $K=\Q(a)$ are obtained from the factorization
$P=\prod g_i$ of the minimal polynomial over $K$ (\code{Factor[P, Extension -> Root[P,k]]}, $0.3$ s for
the examples): each factor $g$ gives the principal subfield $V_g=\{h(\theta): g\mid h(x)-h(\theta)\}$,
computed as a null space, and every subfield is an intersection of principal subfields
\cite{SvH,r5,r9}.

For completeness, the field property and the intersection assertion have a short explanation.  A
polynomial $g\mid P$ imposes $h(\beta)=h(\theta)$ for its roots $\beta$, with $h$ taken modulo $P$;
these equalities are preserved by addition, multiplication and inversion of a nonzero element, so
$V_g$ is a subfield of $K$.  If $E\subseteq K$, the roots of the minimal polynomial of $\theta$ over
$E$ are exactly the images of $\theta$ under embeddings of $K$ fixing $E$.  An element $h(\theta)$
has the same value under all those embeddings if and only if it belongs to $E$, by separability.
Factoring that minimal polynomial over $K$ therefore expresses $E$ as the intersection of the
corresponding $V_{g_i}$.  Closing the principal-field list under intersections therefore gives all
subfields, with an exact linear-algebra construction.

In $K$ the additive criterion is complete for sums inside $K$, and the two-factor
criterion with $t=1$ is complete for two factors inside $K$.  If $K$ is Galois (all $g_i$ linear),
the additive search is globally complete and the product search becomes globally complete for two
factors by allowing all the exponents of Theorem~\ref{thm:two}.  The option restricting components to
$K$ must still enforce $t=1$, even for Galois $K$.  The implementations try this path first and stop when it reaches a lower bound,
which is what happens for both examples in \eqref{eq:examples}: the answers \eqref{eq:answers} are found
and certified in about a second, without constructing the splitting field.

\subsection{Normalization of the output}

Linear algebra returns arbitrary elements of the solution spaces.  Summands are trace-centred (the
rational mean of the conjugates, $\Tr(y)/[L:\Q]$, is subtracted and collected into a rational offset);
this recovers the trace-zero cubics of \eqref{eq:answers}.  If a term-count bound is requested, the
offset is absorbed into one summand instead of adding a term; a rational translation does not change
the field generated by that summand.  For products the null vector is taken from
an LLL-reduced basis of the integer null space, the factor $u$ is rescaled by the $p$-adically minimal
rational $q$ making $qu$ an algebraic integer (for each prime, the least exponent $e$ with
$v_p(c_i)+(d-i)e\ge0$ for all coefficients $c_i$ of the monic minimal polynomial), followed by a small
search for the scaling of least coefficient height; the cofactor is divided accordingly.  This turns
$\tfrac1{4}\cdot\Root[8+4\#+\#^3\&,1]\cdot\Root[-2+\ldots]$ into the factors of \eqref{eq:answers}.

\subsection{Interface}

\begin{lstlisting}
Get["RootDecomposition.wl"];
RootProductDecomposition[a]      RootProductDecomposition[a, d]
RootSumDecomposition[a]          RootSumDecomposition[a, d]
RootGaloisData[a]                RootDecompositionLowerBound[a]
RootBoundedDecomposition[a, Plus|Times, d, height, factors]
RootDecompositionVerify[a, terms, Plus|Times]
\end{lstlisting}
An association reports the components in \code{"Terms"}, their \code{"Degrees"} and
\code{"MaximumDegree"}, and a \code{"LowerBound"}.  The certificate fields are \code{"Optimal"}
(global optimality, with the qualifications below), \code{"ScopeOptimal"} (optimal within the
requested class) and \code{"TwoFactorOptimal"}.  Further fields record the \code{"NormExponent"},
\code{"Verified"} status, \code{"Method"} and an inactive \code{"Expression"}.

\begin{sloppypar}
To keep components inside $\Q(a)$, set \code{"Scope"} to \code{"InputField"}.
Set \code{"Engine"} to \code{"SplittingField"} to force the Galois route.
Setting \code{"Coefficients"} to \code{"GaussianRationals"} enables unrestricted Gaussian sums, while
\code{"MaxTerms"} restricts the number of ordinary summands.  Combining Gaussian coefficients with
a finite term cap is rejected because the unrestricted linear-space calculation does not solve that
problem.  Likewise, an explicit degree cap must either be met or produce a failure, including when a
trivial representation is returned.
\code{"MaxFactors" -> 2} restricts to the complete two-factor algorithm; \code{"MaxGroupOrder"} and
\code{"WorkingPrecision"} bound the Galois computation.
\end{sloppypar}

Python's two search functions are \code{sum\_decomposition} and
\code{product\_decomposition}, both in \code{rootdecomp.py}.  The module also exposes
\code{galois\_data} and \code{lower\_bound}, the verifier \code{verify\_exact}, and the input parser
\code{parse\_wolfram\_root}.  The ordinary decomposition scopes have the same
meaning; Gaussian coefficients are a Wolfram-only feature, and the root-index portability
qualification in the conventions applies to non-real branches.

\subsection{Measurements}

Table~\ref{tab:timing} records the original development measurements (Wolfram 15.0.1, Python 3.14
with python-flint 0.8).  These historical timings are not fresh measurements of the revised code.
In that run the Galois route of the Python port was faster than the Wolfram kernel because
exact factorization of the resolvents in FLINT is cheaper than \code{RootReduce} of the partial primitive
elements, and its rational linear algebra runs in FLINT rather than in the interpreter.  The fast path
in $\Q(a)$ is the best choice for both examples.

\begin{table}[htbp]
\centering\small
\begin{tabular}{@{}lrr@{}}
\toprule
Task & Wolfram & Python\\
\midrule
Product example, fast path in $\Q(a)$, certified optimal & 0.2 s & 1.9 s\\
Sum example, fast path in $\Q(a)$, certified optimal & 1.3 s & 0.4 s\\
Splitting-field data of the degree-9 examples ($|G|=36$) & 18--21 s & 3--5 s\\
$\Dp(\alpha_\times)=6$ with the splitting-field data cached & 0.2 s & 0.15 s\\
$\sqrt{(1+\sqrt2)(1+\sqrt3)}$, $|G|=16$, two-factor with $t=2$ & 4 s & 0.4--2 s\\
$x^4-x-1$ pair sum, $|G|=24$ & 3 s & 0.5 s\\
Product search for $\alpha_+$ (two-factor, tensor families, recursion) & 30 s & 110 s\\
\code{ToNumberField[roots, All]} for the degree-9 examples & 14 min & --\\
\bottomrule
\end{tabular}
\caption{Historical wall-clock times (Wolfram 15.0.1; Python 3.14 with python-flint 0.8 and SymPy 1.14).
Current validation and benchmark results are recorded separately in the project documentation.}\label{tab:timing}
\end{table}

\subsection{Limits and scope of the certificates}

The implementations of the complete field-based searches need certified field data; the resolvent method is practical for groups of order
up to a few hundred (the resolvent at level $k$ has degree up to $n\,|G|$).  For larger groups the fast
path still returns decompositions inside $\Q(a)$ and the bounded search still finds small-height
factors, but an unsuccessful bounded search gives no global negative statement.

Three distinct assertions need separate evidence: an identity holds for the selected branches; the
reported component degrees are correct; and no better representation exists in the stated scope.
The Python Galois construction uses ball arithmetic and exact factorization, and rational linear
algebra is then exact.  Its independent \code{verify\_exact} reconstructs the sum or product by exact
polynomial arithmetic with certified root selection; \code{verify\_numeric} is only an enclosure
compatibility diagnostic.  Polynomial divisibility without branch selection, or overlapping
numerical balls alone, would not prove a particular identity.

The Wolfram splitting-field route relies on exact minimal polynomials, numerical root matching with
generous margins, and group-closure and automorphism-consistency checks.  These are useful checks but
do not replace a proof of every numerical reconstruction step.  An exactly verified identity attaining
the independent degree or Frobenius lower bound certifies global optimality without that reconstruction.
An optimality flag relying on exhaustive Wolfram splitting-field enumeration instead is conditional
on the reconstructed group and field data.  Exact verification of the returned identity alone does not
remove this qualification.  A returned \code{"Optimal" -> False} means \emph{not established}, not
\emph{suboptimal}; \code{"TwoFactorOptimal"} concerns the two-factor problem, and a request limiting
the number of terms concerns that restricted class rather than the unrestricted minimum.

\section{Conclusion}

The additive problem is completely solved by a finite algorithm, and the two-factor multiplicative
problem likewise; the key steps are trace descent to the splitting field and norm descent with an
explicit radical exponent, both of which reduce to rational linear algebra on fixed fields.  The
arbitrary-length multiplicative problem is different in kind: its answer depends on arithmetic
invisible to the Galois group, and the present algorithms do not decide it in general.  For the two numbers of the question the
answer is unambiguous --- the cubic decompositions \eqref{eq:answers} are globally optimal --- and the
same machinery shows that the product root needs sextic summands, that optimal factors can live outside
the splitting field, and that flat sums of three quadratic numbers can be inaccessible to any binary
splitting, even with Gaussian coefficients.

\begin{thebibliography}{99}\raggedright\sloppy
\bibitem{question} V. Reshetnikov, \emph{Factor a polynomial Root into Roots of smallest possible
  degree}, Mathematica Stack Exchange, question 105933 (2016), with an answer by yarchik (2019).
  \url{https://mathematica.stackexchange.com/q/105933/7288}.
\bibitem{WolframRoot} Wolfram Research, \emph{Root}, Wolfram Language documentation,
  details on root indexing and the \code{ExactRootIsolation} option.
  \url{https://reference.wolfram.com/language/ref/Root.html}.
\bibitem{r1} Report 1, \emph{Decomposing Algebraic Numbers into Low-Degree Roots} (2026),
  \code{reports/report-01} in the repository.
\bibitem{r2} Report 2, \emph{Decomposing Algebraic Numbers into Low-Degree Roots: Certified Mathematica
  searches, global degree bounds, and exact normal-field algorithms} (2026), \code{reports/report-02}.
\bibitem{r3} Report 3, \emph{Decomposing Polynomial Roots into Low-Degree Sums and Products} (2026),
  \code{reports/report-03}.
\bibitem{r4} Report 4, \emph{Decomposing Algebraic Numbers into Low-Degree Roots: Exact Mathematica
  methods, global minimality certificates, and the distinction between sums and products} (2026),
  \code{reports/report-04}.
\bibitem{r5} Report 5, \emph{Decomposing algebraic numbers into low-degree Roots} (2026),
  \code{reports/report-05}.
\bibitem{r6} Report 6, \emph{Decomposing Algebraic Roots into Low-Degree Sums and Products} (2026),
  \code{reports/report-06}.
\bibitem{r7} Report 7, \emph{Decomposing Algebraic Roots into Low-Degree Sums and Products: Mathematica
  algorithms, optimality certificates, and the role of the splitting field} (2026),
  \code{reports/report-07}.
\bibitem{r8} Report 8, \emph{Decomposing Algebraic Roots Into Sums and Products of Minimum Degree}
  (2026), \code{reports/report-08}.
\bibitem{r9} Report 9, \emph{Decomposing Algebraic Numbers into Low-Degree Sums and Products: Exact
  algorithms, global degree bounds, and a Wolfram Language reference implementation} (2026),
  \code{reports/report-09}.
\bibitem{Altamirano} C. Altamirano, \emph{A Problem in Algebraic Number Theory}, MIT UROP+ final
  paper (2017), as cited in \cite{r1,r3,r7}.
\bibitem{Samuels} C. L. Samuels, \emph{The infimum in the metric Mahler measure}, arXiv:1408.4165
  (2014).
\bibitem{BFSS} A. Bostan, P. Flajolet, B. Salvy, \'E. Schost, \emph{Fast computation of special
  resultants}, J. Symbolic Comput. 41 (2006), 1--29.
\bibitem{Cohen} H. Cohen, \emph{A Course in Computational Algebraic Number Theory}, Graduate Texts in
  Mathematics 138, Springer, 1993.
\bibitem{SvH} J. Szutkoski, M. van Hoeij, \emph{The complexity of computing all subfields of an
  algebraic number field}, J. Symbolic Comput. 93 (2019), 161--182; arXiv:1606.01140.
\bibitem{vHKN} M. van Hoeij, J. Kl\"uners, A. Novocin, \emph{Generating subfields}, J. Symbolic
  Comput. 52 (2013), 17--34.
\bibitem{Stauduhar} R. P. Stauduhar, \emph{The determination of Galois groups}, Math. Comp. 27
  (1973), 981--996.
\bibitem{FK} C. Fieker, J. Kl\"uners, \emph{Computation of Galois groups of rational polynomials},
  LMS J. Comput. Math. 17 (2014), 141--158.
\bibitem{Arb} F. Johansson, \emph{Arb: efficient arbitrary-precision midpoint-radius interval
  arithmetic}, IEEE Trans. Comput. 66 (2017), 1281--1292.
\bibitem{FLINT} W. Hart, \emph{Fast Library for Number Theory: an introduction}, in: Mathematical
  Software -- ICMS 2010, Lecture Notes in Computer Science 6327, Springer, 2010, 88--91;
  python-flint, \url{https://github.com/flintlib/python-flint}.
\bibitem{Milne} J. S. Milne, \emph{Fields and Galois Theory}, course notes, version 5.10 (2022),
  \url{https://www.jmilne.org/math/CourseNotes/FT.pdf}.
\end{thebibliography}

\end{document}
